% --- GLOBAL VARIABLES ---
\newcommand{\baseradius}{4}
\newcommand{\tubeborder}{.8}
% \tube{<comma separated list of coordinate/node names, e.g. A,B,C,A,D,E,C>}{<nesting depth>}
%
% Draws a tube (blue outline, white interior) of radius \baseradius + 2*(depth) pt
% around the closed polygonal path through the listed points, filling its interior
% (as the original version did via \filldraw ... -- cycle).
%
% Improvements over the original:
%  * The path is assembled *outside* the path parser (the original \ifnum inside
%    \filldraw silently lost the initial moveto, so nothing between the vertices
%    was ever drawn).
%  * All radii and line widths are measured in canvas pt while the coordinate
%    transformation is reset (\pgftransformreset). Node coordinates are stored
%    absolutely, so vertex disks stay perfect circles and their radius equals the
%    half-width of the stroked segments exactly. Hence the boundary is a single
%    C^1 curve (disks tangent to the straight strips), and this remains true when
%    the picture is scaled -- also separately in x and y. Only positions scale;
%    tube thickness and ring spacing stay constant, consistent across a document.
\newcommand{\tube}[2]{%
\begin{scope}[on background layer]
    \pgfmathsetmacro{\currentradius}{\baseradius + (#2 * 2)}%
    % assemble "(v1) -- (v2) -- ..." and count the vertices
    \gdef\tubepath{}\gdef\tubecount{0}%
    \foreach \v [count=\c] in {#1}{%
        \ifnum\c=1 \xdef\tubepath{(\v.center)}\else\xdef\tubepath{\tubepath -- (\v.center)}\fi
        \xdef\tubecount{\c}%
    }%
    \pgftransformreset % radii/line widths below are true canvas lengths
    \foreach \v in {#1}{\fill[Blue] (\v) circle (\currentradius pt);}%
    \ifnum\tubecount>1
        \filldraw[Blue, line width = 2*\currentradius pt, line join=round]
            \tubepath -- cycle;
    \fi
    \foreach \v in {#1}{\fill[white] (\v) circle (\currentradius pt - \tubeborder pt);}%
    \ifnum\tubecount>1
        \filldraw[white, line width = 2*\currentradius pt - 2*\tubeborder pt, line join=round]
            \tubepath -- cycle;
    \fi
\end{scope}
}
% a copy of the path P4 with vertices (#1-1)...(#1-4), centered at the current origin
\newcommand{\pathgraph}[1]{
    \foreach\x in {1,2,3,4}{
        \node[fill, circle, inner sep=1.5pt] (#1-\x) at ({(\x-2.5)*0.5},0) {};
    }
    \draw (#1-1) -- (#1-4);
}

\begin{tikzpicture}
    % invisible boxes at the vertices of the associahedron (used to clip the edges)
    \node[rectangle, minimum width=2.26cm, minimum height=0.88cm, inner sep=0pt] (box-1-12-123) at (-2.64,0.17) {};
    \node[rectangle, minimum width=2.26cm, minimum height=0.88cm, inner sep=0pt] (box-1-12-4) at (-0.83,-0.83) {};
    \node[rectangle, minimum width=2.26cm, minimum height=0.88cm, inner sep=0pt] (box-1-123-3) at (-1.50,1.50) {};
    \node[rectangle, minimum width=2.26cm, minimum height=0.88cm, inner sep=0pt] (box-1-3-34) at (-0.17,2.64) {};
    \node[rectangle, minimum width=2.26cm, minimum height=0.88cm, inner sep=0pt] (box-1-34-4) at (0.83,0.83) {};
    \node[rectangle, minimum width=2.26cm, minimum height=0.88cm, inner sep=0pt] (box-12-123-2) at (-2.34,-2.34) {};
    \node[rectangle, minimum width=2.26cm, minimum height=0.88cm, inner sep=0pt] (box-12-2-4) at (0.17,-2.64) {};
    \node[rectangle, minimum width=2.26cm, minimum height=0.88cm, inner sep=0pt] (box-123-2-23) at (-3.90,-3.90) {};
    \node[rectangle, minimum width=2.26cm, minimum height=0.88cm, inner sep=0pt] (box-123-23-3) at (-3.79,3.79) {};
    \node[rectangle, minimum width=2.26cm, minimum height=0.88cm, inner sep=0pt] (box-2-23-234) at (3.79,-3.79) {};
    \node[rectangle, minimum width=2.26cm, minimum height=0.88cm, inner sep=0pt] (box-2-234-4) at (1.50,-1.50) {};
    \node[rectangle, minimum width=2.26cm, minimum height=0.88cm, inner sep=0pt] (box-23-234-3) at (3.90,3.90) {};
    \node[rectangle, minimum width=2.26cm, minimum height=0.88cm, inner sep=0pt] (box-234-3-34) at (2.34,2.34) {};
    \node[rectangle, minimum width=2.26cm, minimum height=0.88cm, inner sep=0pt] (box-234-34-4) at (2.64,-0.17) {};
    % edges of the associahedron
    \draw (box-1-12-123) -- (box-1-12-4);
    \draw (box-1-12-123) -- (box-1-123-3);
    \draw (box-1-12-123) -- (box-12-123-2);
    \draw (box-1-12-4) -- (box-1-34-4);
    \draw (box-1-12-4) -- (box-12-2-4);
    \draw (box-1-123-3) -- (box-1-3-34);
    \draw (box-1-123-3) -- (box-123-23-3);
    \draw (box-1-3-34) -- (box-1-34-4);
    \draw (box-1-3-34) -- (box-234-3-34);
    \draw (box-1-34-4) -- (box-234-34-4);
    \draw (box-12-123-2) -- (box-12-2-4);
    \draw (box-12-123-2) -- (box-123-2-23);
    \draw (box-12-2-4) -- (box-2-234-4);
    \draw (box-123-2-23) -- (box-123-23-3);
    \draw (box-123-2-23) -- (box-2-23-234);
    \draw (box-123-23-3) -- (box-23-234-3);
    \draw (box-2-23-234) -- (box-2-234-4);
    \draw (box-2-23-234) -- (box-23-234-3);
    \draw (box-2-234-4) -- (box-234-34-4);
    \draw (box-23-234-3) -- (box-234-3-34);
    \draw (box-234-3-34) -- (box-234-34-4);
    % vertices of the associahedron = maximal tubings of the path
    \begin{scope}[shift={(-2.64,0.17)}]
        \pathgraph{1-12-123}
        \tube{1-12-123-1,1-12-123-3}{3}
        \tube{1-12-123-1,1-12-123-2}{2}
        \tube{1-12-123-1}{1}
    \end{scope}
    \begin{scope}[shift={(-0.83,-0.83)}]
        \pathgraph{1-12-4}
        \tube{1-12-4-1,1-12-4-2}{2}
        \tube{1-12-4-1}{1}
        \tube{1-12-4-4}{1}
    \end{scope}
    \begin{scope}[shift={(-1.50,1.50)}]
        \pathgraph{1-123-3}
        \tube{1-123-3-1,1-123-3-3}{2}
        \tube{1-123-3-1}{1}
        \tube{1-123-3-3}{1}
    \end{scope}
    \begin{scope}[shift={(-0.17,2.64)}]
        \pathgraph{1-3-34}
        \tube{1-3-34-3,1-3-34-4}{2}
        \tube{1-3-34-1}{1}
        \tube{1-3-34-3}{1}
    \end{scope}
    \begin{scope}[shift={(0.83,0.83)}]
        \pathgraph{1-34-4}
        \tube{1-34-4-3,1-34-4-4}{2}
        \tube{1-34-4-1}{1}
        \tube{1-34-4-4}{1}
    \end{scope}
    \begin{scope}[shift={(-2.34,-2.34)}]
        \pathgraph{12-123-2}
        \tube{12-123-2-1,12-123-2-3}{3}
        \tube{12-123-2-1,12-123-2-2}{2}
        \tube{12-123-2-2}{1}
    \end{scope}
    \begin{scope}[shift={(0.17,-2.64)}]
        \pathgraph{12-2-4}
        \tube{12-2-4-1,12-2-4-2}{2}
        \tube{12-2-4-2}{1}
        \tube{12-2-4-4}{1}
    \end{scope}
    \begin{scope}[shift={(-3.90,-3.90)}]
        \pathgraph{123-2-23}
        \tube{123-2-23-1,123-2-23-3}{3}
        \tube{123-2-23-2,123-2-23-3}{2}
        \tube{123-2-23-2}{1}
    \end{scope}
    \begin{scope}[shift={(-3.79,3.79)}]
        \pathgraph{123-23-3}
        \tube{123-23-3-1,123-23-3-3}{3}
        \tube{123-23-3-2,123-23-3-3}{2}
        \tube{123-23-3-3}{1}
    \end{scope}
    \begin{scope}[shift={(3.79,-3.79)}]
        \pathgraph{2-23-234}
        \tube{2-23-234-2,2-23-234-4}{3}
        \tube{2-23-234-2,2-23-234-3}{2}
        \tube{2-23-234-2}{1}
    \end{scope}
    \begin{scope}[shift={(1.50,-1.50)}]
        \pathgraph{2-234-4}
        \tube{2-234-4-2,2-234-4-4}{2}
        \tube{2-234-4-2}{1}
        \tube{2-234-4-4}{1}
    \end{scope}
    \begin{scope}[shift={(3.90,3.90)}]
        \pathgraph{23-234-3}
        \tube{23-234-3-2,23-234-3-4}{3}
        \tube{23-234-3-2,23-234-3-3}{2}
        \tube{23-234-3-3}{1}
    \end{scope}
    \begin{scope}[shift={(2.34,2.34)}]
        \pathgraph{234-3-34}
        \tube{234-3-34-2,234-3-34-4}{3}
        \tube{234-3-34-3,234-3-34-4}{2}
        \tube{234-3-34-3}{1}
    \end{scope}
    \begin{scope}[shift={(2.64,-0.17)}]
        \pathgraph{234-34-4}
        \tube{234-34-4-2,234-34-4-4}{3}
        \tube{234-34-4-3,234-34-4-4}{2}
        \tube{234-34-4-4}{1}
    \end{scope}
\end{tikzpicture}
